% Generated draft: requires mathematical and source review.
\documentclass[11pt]{article}
\usepackage[margin=25mm]{geometry}
\usepackage{fontspec}
\setmainfont{Latin Modern Roman}
\usepackage{amsmath,amssymb,mathtools,unicode-math}
\setmathfont{Latin Modern Math}
\usepackage{enumitem,needspace,xurl,hyperref}
\setlength{\parindent}{0pt}
\setlength{\parskip}{5pt}
\setlength{\emergencystretch}{3em}
\newcommand{\N}{\mathbb{N}}
\newcommand{\Z}{\mathbb{Z}}
\newcommand{\Q}{\mathbb{Q}}
\newcommand{\R}{\mathbb{R}}
\newcommand{\C}{\mathbb{C}}
\newcommand{\F}{\mathbb{F}}
\newcommand{\A}{\mathbb{A}}
\newcommand{\PP}{\mathbb{P}}
\newcommand{\E}{\mathbb{E}}
\newcommand{\eps}{\varepsilon}
\newcommand{\dd}{\,\mathrm{d}}
\begin{document}
\section*{UnsolvedMath ID: 800049}
\textbf{Source problem code:} catalogue entry
\subsection{Definitions and mathematical statement}
The supplied source uses the term \emph{pargoid of transformations} without defining its partial operation, its isomorphisms, or its endomorphisms. Consequently, the following is a reconstruction schema; a fully specified problem requires those conventions and the intended classes of objects.

For a specified class \(\mathcal C\) of mathematical structures, let \(T\) assign to each \(A\in\mathcal C\) a specified pargoid \(T(A)\) of transformations associated with \(A\). The first task is to determine classes \(\mathcal D\subseteq\mathcal C\) for which
\[
\forall A,B\in\mathcal D,\qquad
T(A)\cong T(B)\ \Longrightarrow\ A\cong B.
\]
Here the left-hand isomorphism is an isomorphism of pargoids and the right-hand isomorphism is an isomorphism of the original structures. The source requests this for the main classes of mathematical structures and natural pargoids of transformations, but does not specify which classes or assignments are intended.

For the second task, fix an intended comparison class \(\mathcal K\) of pargoids of transformations. Once the relevant morphisms have been specified, write \(\operatorname{End}(P)\) for the semigroup of endomorphisms of \(P\), with composition, and \(\operatorname{Aut}(P)\) for its automorphism group. Find pargoids \(P\in\mathcal K\) satisfying, respectively,
\[
\forall Q\in\mathcal K,\qquad
\operatorname{End}(P)\cong\operatorname{End}(Q)
\ \Longrightarrow\ P\cong Q,
\]
or
\[
\forall Q\in\mathcal K,\qquad
\operatorname{Aut}(P)\cong\operatorname{Aut}(Q)
\ \Longrightarrow\ P\cong Q.
\]
The invariant isomorphisms are semigroup and group isomorphisms, respectively. These are separate reconstruction questions. The source does not specify the comparison class \(\mathcal K\).

As background, for a partially ordered set \((X,\leq_X)\), a closure operator in the usual sense is a map \(c:X\to X\) such that, for all \(x,y\in X\),
\[
x\leq_X c(x),\qquad c(c(x))=c(x),\qquad
x\leq_X y\ \Longrightarrow\ c(x)\leq_X c(y).
\]
The source attributes to E.\,S. Lyapin the result that partially ordered sets are determined up to isomorphism by the pargoids of their closure operators, citing [ESL2], [ESL3], and [ESL4]. Symbolically, writing \(\mathsf{Cl}(X)\) for the intended pargoid, the reported result is
\[
\mathsf{Cl}(X)\cong\mathsf{Cl}(Y)
\ \Longrightarrow\ (X,\leq_X)\cong(Y,\leq_Y).
\]
The partial operation on \(\mathsf{Cl}(X)\) is not supplied, so this reported result does not itself define the required pargoid construction. It is recorded as attributed background, without independent verification.
\subsection{Short English statement}
Determine which classes of mathematical structures can be reconstructed, up to isomorphism, from specified natural pargoids of transformations. Also find transformation pargoids reconstructible from their endomorphism semigroups or, separately, their automorphism groups. The source leaves the pargoid conventions, natural constructions, and comparison classes unspecified. It attributes a closure-operator reconstruction result for partially ordered sets to E.\,S. Lyapin.
\subsection{Sources}
\begin{itemize}
\item Local UnsolvedMath registry: permanent ID \(800049\), exact source problem code \texttt{catalogue entry}, title \emph{The Lyapin's Notebook --- Problem 51}. The supplied registry labels the problem \texttt{open}; this is supplied status metadata, not an independently verified literature assessment.
\item Bershadsky, Kublanovsky, and Mashevitzky, \emph{The Lyapin's Notebook: Unsolved Problems in Semigroup Theory} (2026), Problem 51. Supplied source URL: \url{https://arxiv.org/abs/2604.04763}. The supplied provenance records access on 2026-10-16 and extraction by \texttt{arxiv\_latex}; neither access nor source contents were independently verified for this draft.
\item The problem statement attributes the questions to E.\,S. Lyapin. The registry's \texttt{proposed\_by} field lists Bershadsky, Kublanovsky, and Mashevitzky, and its \texttt{proposed\_year} field is 2026; these metadata do not establish the original date of Lyapin's questions.
\item The supplied statement cites [ESL2], [ESL3], and [ESL4] for Lyapin's reported reconstruction result concerning partially ordered sets and closure operators. Bibliographic details for these three references are not supplied.
\end{itemize}
\end{document}
